
\section{Condensed Cubesets}\label{sec:condensed-cubeset}

As of now we have only considered the category $\cSet$ of \enquote{discrete} cubesets without adhering to a topology
on the collections of $n$-cubes.
However, a great deal of the theory of nilspaces deals with \emph{topological} cubesets,
called \emph{cubespaces}, especially compact cubespaces, in order to appropriately treat Host-Kra theory.
After all, dynamical systems should always preserve a topological or measure theoretical structure!
Therefore, in \cite{Gutman2020}, the theory of nilspaces is treated topologically from the beginning,  making the arguments more involved,
and in \cite{Candela2017a}, the arguments from \cite{Candela2017} are revisited again to show that they are consistent with the topological aspects.
In both cases, every construction has to be checked for continuity of maps,
or the Hausdorff property of quotients, etc.
Moreover, the arguments usually are only considered for compact second-countable Hausdorff spaces (= metrizable compact spaces).

In this section we go a different route to install a topology on a cubeset
that is more global and doesn't adhere to the concrete constructions made in the previous sections.
One of the main differences to classical approaches is that we don't work with topological spaces,
but rather a replacement thereof: condensed sets.
The category of condensed sets $\CondSet$ is categorically much better behaved
than the category of topological spaces as it behaves just like a (Grothendieck) topos
(see also Remark~\ref{rem:size-issues}) on the one hand,
and like a category of algebras of an algebraic theory on the other hand.
Nevertheless, the category of condensed sets contains the category of compactly generated weak Hausdorff spaces as a full subcategory
and thus there is no loss of information when one wants to recover (most) topological spaces.
The good categorical properties, especially its definition as a category of sheaves on a coherent site,
allow one to apply the theory of presentable ($\infty$-)categories.
We achieve the topological enrichment by tensoring the category $\cSet$ of cubesets
with the category of condensed sets $\CondSet$ in $\PrL$, the $\infty$-category of presentable $\infty$-categories.


\subsection{Tensor Products of Presentable Categories}

We begin by recalling the tensor product of presentable $(\infty,1)$-categories from \cite[Section 4.8.1]{Lurie2017},
see also \cite[Section 0713]{Kerodon2026}.
In the following, $\infty$-category will always mean $(\infty,1)$-category.
Denote by $\PrL$ the $\infty$-category consisting of presentable categories
whose morphisms are functors $\cC\to\cD$ preserving small colimits, see \cite[Section 5.5]{Lurie2009}.
By the adjoint functor theorem, any colimit preserving functor $\cC\to\cD$ between presentable categories is left adjoint.
We denote the full sub-$\infty$-category of $\Fun(\cC,\cD)$ consisting of colimit preserving functors by $\Fun^\uL(\cC,\cD)$.

The category $\PrL$ has a symmetric monoidal tensor product $\otimes\colon\PrL\times\PrL\to\PrL$.
It represents separately colimit preserving functors, i.e., there is an equivalence of $\infty$-categories
\[
    \Fun^\uL(\cC\otimes\cD,\cE) \simeq \Fun_{\mathrm{colim}\times\mathrm{colim}}(\cC\times\cD,\cE)
\]
where the right hand side denotes the full sub-$\infty$-category of functors $\cC\times\cD\to\cE$
that preserve colimits in both variables separately.
Moreover, the tensor structure is closed with internal Hom object $\Fun^\uL(\cC,\cE)$,
i.e., there is a natural equivalence
\[
    \Fun^\uL(\cC\otimes\cD,\cE) \simeq \Fun^\uL(\cD,\Fun^\uL(\cC,\cE)).
\]
We denote the tensor unit of $\PrL$ by $\Ani$ (which is the free cocompletion of a point).
To give an explicit description of the tensor product we need the following definition.

\begin{definition}\label{def:cont-fun}
Let $\cK$ be a family of small $\infty$-categories and $S,\cD$ $\infty$-categories with all limits of forms in $\cK$.
Then we define $\Fun^\cK(S,\cD)$ to be the full sub-$\infty$-category of $\Fun(S,\cD)$ spanned by functors preserving $\cK$-indexed limits.
\end{definition}

\begin{example}\label{ex:cont-fun-cats}
\begin{enumerate}[(i)]
    \item If $\cK$ consists of all small $\infty$-categories and $S=\cC^\op$ with $\cC$ a presentable $\infty$-category,
            then we have that $\Fun^\cK(\cC^\op,\cD)$ consists of all limit preserving functors $\cC^\op\to\cD$
            which equivalently are right adjoint functors, and we denote this $\infty$-category by $\Fun^\uR(\cC^\op,\cD)$.
    \item If $\cK$ is empty, then $\Fun^\cK(S,\cD) = \Fun(S,\cD)$.
    \item If $\cK$ consists of all finite $\infty$-categories, $S$ is small and $\cD$ a presentable $\infty$-category,
            then we denote $\Fun^\cK(S,\cD)$ by $\Fun^\mathrm{lex}(S,\cD)$ (= left exact functors)
            which again is a presentable $\infty$-category.
    \item If $\cK$ consists of all finite sets, $S$ is small and $\cD$ a presentable $\infty$-category,
            we write $\Fun^\times(S,\cD)$ for $\Fun^\cK(S,\cD)$ which is the presentable $\infty$-category of finite product preserving functors.            
\end{enumerate}
\end{example}

\begin{proposition}\label{prop:pres-tens-prod}
Let $\cC$ and $\cD$ be presentable $\infty$-categories.
Then there is a canonical equivalence
\[
    \cC\otimes\cD \simeq \Fun^\uR(\cC^\op,\cD) \simeq \Fun^\uL(\cC,\cD^\op)^\op
\]
\end{proposition}

\begin{proof}
This is \cite[Proposition 4.8.1.17]{Lurie2017}.
\end{proof}


\begin{corollary}\label{cor:pres-tens-prod}
Let $\cK$ be a collection of small $\infty$-categories and $S$ a small $\infty$-category with limits of form in $\cK$.
For $\cC = \Fun^\cK(S,\Ani)$ we have that
\[
    \cC\otimes\cD \simeq \Fun^\cK(S,\cD).
\]
\end{corollary}

\begin{proof}
We have that
\[
    \cC\otimes\cD\simeq\Fun^\uR(\cC^\op,\cD)\simeq\Fun^\uL(\Fun^\cK((S^\op)^\op,\Ani),\cD^\op)^\op\simeq\Fun_\cK(S^\op,\cD^\op)^\op\simeq\Fun^\cK(S,\cD)
\]
where the second to last equivalence follows from \cite[Theorem 5.3.6.2]{Lurie2009}.
\end{proof}

We will also need the following lemma on the functoriality of the tensor product.
In some special cases, this is \cite[Section 2.2]{Haine2025}
but we weren't able to find a complete reference, so we record it here for completeness.

\begin{lemma}\label{lem:pres-tens-prod-fun}
Let $\cK$ be a family of small $\infty$-categories, $S$ a small $\infty$-category and $\cC = \Fun^\cK(S,\Ani)$.
If $f^\ast\colon\cD\to\cD'$ is a map in $\PrL$ that preserves limits of forms in $\cK$,
then the diagram
\begin{center}
\begin{tikzcd}
	{\cC\otimes\cD} & {\cC\otimes\cD'} \\
	{\Fun^\cK(S,\cD)} & {\Fun^\cK(S,\cD')}
	\arrow["{\id\otimes f^\ast}", from=1-1, to=1-2]
	\arrow["\simeq"', from=1-1, to=2-1]
	\arrow["\simeq", from=1-2, to=2-2]
	\arrow["{(f^\ast)\circ -}"', from=2-1, to=2-2]
\end{tikzcd}
\end{center}
is commutative.
\end{lemma}

\begin{proof}
By \cite[Remark 071F]{Kerodon2026}, under the equivalence $\cC\otimes\cD\simeq\Fun^\uR(\cC^\op,\cD)$,
the right adjoint part of $\id\otimes f^\ast$ is given by postcomposition
\[
    f_\ast\circ -\colon\Fun^\uR(\cC^\op,\cD')\to\Fun^\uR(\cC^\op,\cD)
\]
with the right adjoint $f_\ast$ of $f^\ast$.
Thus, under the equivalence $\cC\otimes\cD\simeq\Fun^\cK(S,\cD)$,
the right adjoint of $\id\otimes f^\ast$ is also given by postcomposition with $f_\ast$.
But the left adjoint of postcomposition with $f_\ast$ on the whole functor category is given by
\[
    f^\ast\circ -\colon\Fun(S,\cD)\to\Fun(S,\cD').
\]
Since $f^\ast$ preserves $\cK$-indexed limits, it restricts to a functor
\[
    f^\ast\circ -\colon\Fun^\cK(S,\cD)\to\Fun^\cK(S,\cD')
\]
which then also is the left adjoint to $f_\ast\circ -$.
\end{proof}

Furthermore, the tensor product of presentable categories behaves nicely under truncation.

\begin{example}\label{ex:pres-tens-prod}
\begin{enumerate}[(i)]
    \item If $n\geq -2$, then $\cC\otimes\tau_n\Ani = \tau_n\cC$
        where $\tau_n$ denotes the left adjoint to the inclusion of $n$-truncated objects in $\cC$ and $\Ani$, respectively.
        In particular, $\tau_n\cC\otimes\cD = \cC\otimes\tau_n\cD = \tau_n\cC\otimes\tau_n\cD$.
    \item For $n = 0$ we have that $\tau_0\Ani = \set$ and so $\cC\otimes\set = \tau_0\cC$.
    \item If $\cC$ is a 1-category and $\cD$ any $\infty$-category, then
        $\cC\otimes\cD = \tau_0\cC\otimes\cD = \cC\otimes\tau_0\cD$ which is a 1-category.
    \item The 1-category $\set$ is an idempotent algebra in $\PrL$.
        By the previous we have that
        \[
            \set\otimes\set = \set\otimes\tau_0\Ani = \tau_0\set = \set.
        \]
    \item In particular, the tensor product restricts to 1-categories:
        We have that, if $\cC$ and $\cD$ are presentable 1-categories, then $\cC\otimes_\set\cD = \cC\otimes\cD$,
        and $\set$ is the tensor unit there.
        For this special case see also \cite{Bird1984}.
\end{enumerate}
\end{example}


\subsection{Condensed Cubesets}

To pass from \enquote{discrete} cubesets to \enquote{topological} cubesets we want to apply this tensor product
to the case where $\cC$ is the category of cubesets and $\cD$ is the category of condensed sets.
Let us first recall the category of condensed sets.

We implicitly fix an uncountable strong limit cardinal $\kappa$.

\begin{definition}\label{def:cond}
\begin{enumerate}[(i)]
    \item Denote by $\extr$ the category of extremally disconnected compact Hausdorff spaces (and continuous maps)
            and by $\extr_\kappa$ its full subcategory of spaces with topological weight less than $\kappa$.
    \item We denote the sifted Ind-completion of $\extr_\kappa$ by $\CondAni$, the $\infty$-category of \emph{($\kappa$-)condensed anima}.
            Its category of (homotopy!) discrete objects (0-truncation) is denoted by $\CondSet$,
            the category of \emph{($\kappa$-)condensed sets}.
\end{enumerate}
\end{definition}

Condensed sets have been introduced in \cite{Scholze2026,Scholze2026a,Clausen2026}
with a particular focus on applications in analytic and complex geometry.
For a thorough introduction to the basics of condensed sets see \cite{Bihlmaier2025}.

\begin{proposition}\label{prop:cond-sheaves}
There are equivalences of ($\infty$-)categories
\[
    \CondAni\simeq\Fun^\times(\extr_\kappa^\op,\Ani)\quad\text{and}\quad\CondSet\simeq\Fun^\times(\extr_\kappa^\op,\set).
\]
\end{proposition}

\begin{proof}
This follows from \cite[Proposition 5.5.8.15, Remark 5.5.8.16]{Lurie2009} in the $\Ani$-valued case
and \cite[Corollary 2.8]{Adamek2001} in the $\set$-valued case.
\end{proof}

\begin{remark}\label{rem:size-issues}
At first it might seem arbitrary to fix a strong limit cardinal $\kappa$.
We mainly do so in order to ensure that $\CondAni$ and $\CondSet$ are presentable
since they are ($\infty$-, resp. 1-)sheaf categories on a small site
and thus we can simply apply all aspects of the theory of presentable $\infty$-categories without further ado.
However, one can certainly take the colimit over all $\kappa$ and obtain the \enquote{actual} category of condensed anima/sets,
see also \cite[Remark 1.4]{Scholze2026} and \cite[Theorem 2.4.2]{Bihlmaier2025}.
Then, one has to work with accessible sheaves on $\extr$ and extend the theory of tensor product of presentable categories
to the slightly more general setting to which condensed ($\infty$-)categories belong.
\end{remark}

\begin{remark}\label{rem:cond-sheaves-comparison}
By the comparison lemma, see \cite[Theorem 2.3.10]{Bihlmaier2025}, \cite[Appendix, Corollary 4.3]{MacLane1994} and \cite[Remark 4.4.8]{Kuijper2026},
the $\infty$-category of condensed anima is equivalent to the $\infty$-category of hypersheaves on the site $\prolak$.
Here, $\prolak$ is the category of profinite sets (= Stone spaces) of weight less than $\kappa$,
and covers are given by finite families of jointly surjective maps.
\end{remark}

As usual with a (1-)topos, there are certain finiteness and separation conditions, see \cite[Chapter 2.5]{Bihlmaier2025}.

\begin{definition}\label{def:cond-qcqs}
\begin{enumerate}[(i)]
    \item A cover of a condensed set $X$ is a family $(U_i\to X)_{i\in I}$ such that the map $\coprod_{i\in I}U_i\to X$ is surjective.
    \item A condensed set $X$ is \emph{quasicompact} if every cover has a finite subcover.
    \item A condensed set $X$ is \emph{quasiseparated} if for any two maps $f,g\colon K\to X$ with $K$ quasicompact,
            the pullback $K\times_X K$ is quasicompact.
\end{enumerate}
\end{definition}

A condensed set can be thought of as a replacement for a topological space.
Namely, the assignment
\[
    X \mapsto (E\mapsto C(E,X))
\]
maps a topological space $X$ to a condensed set: its $E$-valued points are simply the continuous functions $E\to X$,
see \cite[Section 2.4]{Bihlmaier2025}.
Thus, this functor is faithful and, if $X$ is ($\kappa$-)compactly generated, then all those functions recover $X$ completely.
This includes locally ($\kappa$-)compact and metrizable spaces.

\begin{proposition}\label{prop:cond-top-compactly-gen}
The functor $\Top\to\CondSet$ is fully faithful on $\kappa$-compactly generated topological spaces
and induces an equivalence of the category $\CHaus_\kappa$ of compact Hausdorff spaces of weight less than $\kappa$
and the full subcategory $\qcqs$ of quasicompact quasiseparated condensed sets.
\end{proposition}

\begin{proof}
See \cite[Corollary 2.3.29]{Bihlmaier2025} for the first, and \cite[Theorem 2.5.46]{Bihlmaier2025} for the second statement.
\end{proof}

\begin{definition}\label{def:cond-cube}
The category of \emph{condensed cubesets} is $\cSet\otimes\CondAni = \cSet\otimes\CondSet$.
We denote this category by $\cCond$.
\end{definition}

\begin{proposition}\label{prop:cond-cube-descr}
Explicitly, the category $\cCond$ of condensed cubesets is canonically equivalent to any of the following categories.
\begin{enumerate}[(a)]
    \item The category of right adjoint functors $\cSet^\op\to\CondSet$,
    \item the category of right adjoint functors $\CondSet^\op\to\cSet$,
    \item the category of functors ${\bbox}^\op\to\CondSet$ \emph{(cubical condensed sets)},
    \item the category of finite product preserving functors $\extr_\kappa^\op\to\cSet$ \emph{(condensed cubesets)},
    \item the category of functors $\extr_\kappa^\op\times{\bbox}^\op\to\set$ that preserve finite products in the first variable.
\end{enumerate}
\end{proposition}

\begin{proof}
The first two follow from Proposition~\ref{prop:pres-tens-prod}.
The equivalence of $(a)$ and $(c)$ follows from Corollary~\ref{cor:pres-tens-prod}, as does the equivalence of $(b)$ and $(d)$.
The equivalence with $(e)$ then follows from currying.
\end{proof}

\begin{remark}
We have already considered a similar construction when talking about cubegroups:
those form the category of functors ${\bbox}^\op\to\Grp$, and so can be described as $\cSet\otimes\Grp$.
Thus, considering topological cubesets (= cubespaces) is essentially the same as working with cubegroups (or any other algebraic object).
The reason for working in condensed sets is that those behave just like any other algebraic theory,
unlike topological spaces, and so admit a description as a sheaf category.
\end{remark}

Recall the inclusion $i\colon\ast\to\bbox$ and the corresponding adjoint functors
$i_!\dashv i^\ast\dashv i_\ast$ from Lemma~\ref{lem:point-cset}.
Also, recall that for a condensed set $X$, the functor given by evaluation at the point $X\mapsto X(\ast)$
admits a left adjoint, given by sending a set $S$ to the sheafification of the constant presheaf
$E\mapsto S$ for $E\in\extr$.
This left adjoint is denoted $S_\ud$ since it agrees with the set $S$
as discrete topological space, viewed as condensed set, see \cite[Lemma 2.4.38]{Bihlmaier2025}.

\begin{lemma}
\begin{enumerate}[(i)]
\item The adjunction $i_!\dashv i^\ast$ induces an adjunction
        $i_!\otimes\id\colon\CondSet\rightleftarrows\cCond\lon i^\ast\otimes\id$.
        Moreover, the diagram
        \begin{center}
        \begin{tikzcd}
            	{\set\otimes\CondSet} & {\cSet\otimes\CondSet} \\
            	\\
            	{\Fun^\times(\extr_\kappa^\op,\set)} & {\Fun^\times(\extr_\kappa^\op,\cSet)}
            	\arrow[""{name=0, anchor=center, inner sep=0}, "{i_!\otimes\id}"', from=1-1, to=1-2]
            	\arrow["\simeq"', from=1-1, to=3-1]
            	\arrow[""{name=1, anchor=center, inner sep=0}, "{i^\ast\otimes\id}"', curve={height=12pt}, from=1-2, to=1-1]
            	\arrow["\simeq", from=1-2, to=3-2]
            	\arrow[""{name=2, anchor=center, inner sep=0}, "{i_!\circ -}"', from=3-1, to=3-2]
            	\arrow[""{name=3, anchor=center, inner sep=0}, "{i^\ast\circ -}"', curve={height=12pt}, from=3-2, to=3-1]
            	\arrow["\dashv"{anchor=center, rotate=84}, draw=none, from=0, to=1]
            	\arrow["\dashv"{anchor=center, rotate=84}, draw=none, from=2, to=3]
        \end{tikzcd}
        \end{center}
        is commutative.
\item The adjunction $(-)_\ud\dashv(-)(\ast)$ induces an adjunction
        $\id\otimes(-)_\ud\colon\cSet\rightleftarrows\cCond\lon \id\otimes(-)(\ast)$.
        Moreover, the diagram
        \begin{center}
        \begin{tikzcd}
	        {\cSet\otimes\set} & {\cSet\otimes\CondSet} \\
            	\\
            	{\Fun({\bbox}^\op,\set)} & {\Fun({\bbox}^\op,\CondSet)}
            	\arrow[""{name=0, anchor=center, inner sep=0}, "{\id\otimes(-)_\ud}"', from=1-1, to=1-2]
            	\arrow["\simeq"', from=1-1, to=3-1]
            	\arrow[""{name=1, anchor=center, inner sep=0}, "{\id\otimes(-)(*)}"', curve={height=12pt}, from=1-2, to=1-1]
            	\arrow["\simeq", from=1-2, to=3-2]
            	\arrow[""{name=2, anchor=center, inner sep=0}, "{(-)_\ud\circ -}"', from=3-1, to=3-2]
            	\arrow[""{name=3, anchor=center, inner sep=0}, "{(-)(*)\circ -}"', curve={height=12pt}, from=3-2, to=3-1]
            	\arrow["\dashv"{anchor=center, rotate=69}, draw=none, from=0, to=1]
            	\arrow["\dashv"{anchor=center, rotate=68}, draw=none, from=2, to=3]
        \end{tikzcd}
        \end{center}
        is commutative.
\end{enumerate}
\end{lemma}

\begin{proof}
That the adjunctions $i_!\dashv i^\ast$ and $(-)_\ud\dashv(-)(\ast)$ upgrade to adjunctions
between the respective tensor products follows from \cite[Section I.1.6.1]{Gaitsgory2017}
where it is shown that the tensor product of presentable $(\infty,1)$-categories
extends to a symmetric monoidal structure on the $(\infty,2)$-category $\PrL$
(where the proof in the unstable case works as in the stable case),
see also \cite[Section 1.2]{Haine2025}.
In particular, the tensor product preserves adjunctions in $\PrL$.

The second claim in (i) follows from Lemma~\ref{lem:pres-tens-prod-fun}
since both $i_!$ and $i^\ast$ preserve limits.
The second claim in (ii) also follows from Lemma~\ref{lem:pres-tens-prod-fun} since $\cK =\emptyset$ in this case.
\end{proof}

\begin{definition}
A condensed cubeset in the essential image of the functor $(-)_\ud\colon\cSet\to\cCond$ is called a \emph{discrete condensed cubeset},
or simply \emph{cubeset}.
\end{definition}

This abuse of notation is justified by the fact that $(-)_\ud$ is fully faithful,
and so its essential image is equivalent to $\cSet$.

Furthermore, there is also a functor $i_!\colon\CondSet\to\cCond$.
This yields, in particular, a Yoneda embedding
\[
    \extr_\kappa\times\bbox\to\cCond,\quad (E,n)\mapsto i_!E\times{\bbox}^n_\ud
\]
since it is given by the composition
\begin{center}
\begin{tikzcd}
	{\extr_\kappa\times\bbox} & {\Fun((\extr_\kappa\times\bbox)^\op,\set)=\Fun(\extr_\kappa^\op,\cSet)} & \cCond.
	\arrow["y", from=1-1, to=1-2]
	\arrow["\Sh", from=1-2, to=1-3]
\end{tikzcd}
\end{center}
To see those two functors to be equal,
note that the Yoneda lemma on the level of presheaves takes the form
\[
    (E,n)\mapsto\left((S,m)\mapsto C(S,E)\times\hom({\bbox}^m,{\bbox}^n)\right)
\]
and since sheafification commutes with products and is computed pointwise on $\bbox$,
we have that
\begin{align*}
       \Sh\left(C(-,E)\times\hom({\bbox}^m,{\bbox}^n)\right)
    &= C(-,E)\times\Sh(\hom({\bbox}^m,{\bbox}^n)) \\
    &= C(-,E)\times C(-,\hom({\bbox}^m,{\bbox}^n)) \\
    &= C(-,E\times\hom({\bbox}^m,{\bbox}^n)).
\end{align*}
The classical Yoneda lemma also gives rise to an enriched Yoneda lemma.
First, note that the Hom-set $\hom(X,Y)$ of two condensed cubesets $X$ and $Y$
naturally has the structure of a condensed set:
its $E$-valued points are given by $\hom(X\times i_!E,Y)$.
This defines a condensed set since $i_!$ preserves coproducts as left adjoint,
so we have that
\[
    X\times i_!(E\sqcup E') = X\times (i_!E\sqcup i_!E') = (X\times i_!E)\sqcup(X\times i_!E')
\]
for $E,E'\in\extr_\kappa$,
where we have used that coproducts are stable under base change in $\CondSet$
and hence by pointwise computation also in $\cCond$.

\begin{lemma}[Enriched Yoneda Lemma]\label{lem:yoneda-enriched}
Let $X$ be a condensed cubeset and $n\in\N_0$.
Then there is an isomorphism of condensed sets
\[
    X(n) \simeq \hom({\bbox}^n_\ud,X).
\]
\end{lemma}

\begin{proof}
We show the claim pointwise on $E\in\extr_\kappa$.
There we have that
\[
    \hom({\bbox}^n_\ud,X)(E) = \hom({\bbox}^n_\ud\times i_!E,X) = X(E,n),
\]
where the last equality follows from the classical Yoneda lemma.
\end{proof}

In the following subsections we will revisit the main notions of cubesets introduced throughout the first chapters
and explain what happens to them after tensoring with $\CondSet$.


\subsection{Concrete Cubesets}

The following proposition allows us to translate properties of discrete cubesets to condensed cubesets.

\begin{proposition}
The concrete cubeset adjunction $(-)^\mathrm{conc}\colon\cSet\rightleftarrows\ccSet$
induces by tensoring with $\CondSet$ an adjunction
$(-)^\mathrm{conc}\colon\cCond\rightleftarrows\cCond^\mathrm{conc}$
with fully faithful right adjoint.
The category $\cCond^\mathrm{conc}$ is canonically identified with the full subcategory of $\cCond$
spanned by objects $X$ for which the map
\[
    X(n) \to \prod_{0\to n} X(0)
\]
is an injection of condensed sets for every $n\in\N$.
Moreover, the left adjoint $(-)^\mathrm{conc}$ is given by
\[
    X^\mathrm{conc}(n) = \im\left(X(n)\to\prod_{0\to n}X(0)\right).
\]
\end{proposition}

\begin{proof}
The existence of the adjunction is clear by functoriality in $\PrL$.
Since both the inclusion $\cSet^\mathrm{conc}\to\cSet$ and its left adjoint $(-)^\mathrm{conc}$ preserve finite products
by Proposition~\ref{prop:concrete-ihom-prod},
under the equivalence $\cSet\otimes\CondSet\simeq\Fun^\times(\extr_\kappa^\op,\cSet)$,
the tensored adjunction is given by pushforward of the original one by Lemma~\ref{lem:pres-tens-prod-fun}.
This implies the rest of the claims.
\end{proof}


\subsection{Fibrant Cubesets}

Since the functor $(-)_\ud\colon\cSet\to\cCond$ preserves colimits,
the corner ${\bbcor}^n_\ud$ is the colimit in $\cCond$ over all inert inclusions ${\bbox}^k_\ud\to{\bbox}^n_\ud$ for $k<n$.

\begin{definition}
A condensed cubeset $X$ is \emph{fibrant} if the map
\[
    X(n) = \hom({\bbox}^n_\ud,X)\to\hom({\bbcor}^n_\ud,X) = \varprojlim_{\substack{{\bbox}^k\to{\bbox}^n\\k<n}}X(k)
\]
is a surjection of condensed sets.
\end{definition}

Equivalently, $X$ is fibrant if for each $E\in\extr_\kappa$, the corresponding cubeset $X(E)$ is fibrant.
Similarly, one defines fibrations of condensed cubesets.

Recall the set of maps
\[
    \bprism\coloneq \{ {{\bbox}^{n}\sqcup_{{\bbox}^{n-1}}{\bbox}^{n}}\to{({\bbox}^{n}\sqcup_{{\bbox}^{n-1}}{\bbox}^{n})\sqcup_{{\bbox}^{n-1}\sqcup{\bbox}^{n-1}}{\bbox}^n} \mid n\in\N \}
\]
from Definition~\ref{def:gluing-property}.

\begin{definition}
A condensed cubeset $X$ has the \emph{gluing property} (\emph{unique gluing property}, \emph{at most unique gluing property}) if the pullback
\[
    f^\ast\colon\hom(B,X)\to\hom(A,X)
\]
is an epimorphism (isomorphism, monomorphism) of condensed sets for every $f\in\bprism$.
\end{definition}

Equivalently, $X$ has the gluing property if for each $E\in\extr_\kappa$, the cubeset $X(E)$ has the gluing property (and likewise for the other properties).

\begin{proposition}
Let $X$ be a condensed cubeset.
Then the following conditions are equivalent.
\begin{enumerate}[(a)]
    \item It has at most unique gluing.
    \item It is concrete.
\end{enumerate}
Moreover, if $X$ is fibrant, then it has the gluing property.
\end{proposition}

\begin{proof}
This follows from testing pointwise against $E\in\extr_\kappa$.
\end{proof}


\subsection{Truncated Cubesets}

Next, we investigate $n$-step condensed cubesets.

\begin{proposition}
Let $n\in\N_0$.
The $n$-truncation $\tau_n\colon\cSet\rightleftarrows\cSet_n$ induces by tensoring with $\CondSet$ an adjunction
$\tau_n\colon\cCond_n\rightleftarrows\cCond$ with fully faithful right adjoint.
The category $\cCond_n$ is canonically identified with the full subcategory of $\cCond$ of $n$-step condensed cubesets $X$:
those for which the map
\[
    \hom({\bbox}^k_\ud,X)\to\hom({\bbcor}^k_\ud,X)
\]
is an isomorphism for every $k>n$.
Moreover, the left adjoint can be computed pointwise on $E\in\extr_\kappa$.
In particular, if $X$ is concrete and fibrant, then the unit $\eta\colon X\to\tau_n X$ is a surjective fibration
and $\tau_n X$ is concrete and fibrant.
\end{proposition}

\begin{proof}
The existence of the $\CondSet$-enriched adjunction follows from functoriality of the tensor product.
Both the inclusion $\cSet_n\to\cSet$ and its left adjoint $\tau_n$ preserve finite products, see Proposition~\ref{prop:n-trunc-ihom-prod}.
Therefore, the adjunction is simply given by pushforward under the equivalence $\cCond\simeq\Fun^\times(\extr_\kappa^\op,\cSet)$
by Lemma~\ref{lem:pres-tens-prod-fun}.
This implies fully faithfulness of the right adjoint and the computation of the left adjoint pointwise on $E\in\extr_\kappa$.
The subcategory $\cCond_n$ is canonically identified with those $X\in\cCond$ such that $X(E)\in\cSet_n$ for every $E\in\extr_\kappa$,
i.e. the map
\[
    X(E)(k) = \hom({\bbox}^k,X(E))\to\hom({\bbcor}^k,X(E)) = (\varprojlim_{j\in\cJ} X(E)(j)) = (\varprojlim_{j\in\cJ} X(j))(E)
\]
is a bijection for every $k>n$.
But this is precisely the condition that the map
\[
    \hom({\bbox}^k_\ud,X)\to\hom({\bbcor}^k_\ud,X)
\]
is an isomorphism of condensed sets for every $k>n$ by the enriched Yoneda lemma~\ref{lem:yoneda-enriched}.
The last claim follows by the pointwise statement on $\extr_\kappa$ from Proposition~\ref{prop:trunc_concfib}.
\end{proof}

\begin{proposition}\label{prop:cond-cube-trunc-qcqs}
Let $X$ be a concrete fibrant condensed cubeset.
If $X$ is (pointwise) qcqs, then so is $\tau_n X$.
\end{proposition}

\begin{proof}
Since the unit $\eta\colon X\to\tau_n X$ is surjective, $\tau_n X$ is quasicompact.
For quasiseparatedness it suffices to see that $\tau_n X(0)$ is quasiseparated:
since $X$ is concrete, the map $X(n)\to\prod_{0\to n}X(0)$ is injective
and subobjects and limits of quasiseparated condensed sets are quasiseparated by \cite[Lemma 2.5.47]{Bihlmaier2025}.
Let $Y = X(0)$ and $Z = \tau_nX(0)$.
Thus it suffices to see that $Z$ is quasiseparated,
which is equivalent to the kernel pair $Y\times_Z Y$ being a closed subobject of $Y\times Y$,
i.e. the inclusion $Y\times_Z Y\to Y\times Y$ is quasicompact.
Since $Y\times Y$ is compact Hausdorff, we need to see that $Y\times_Z Y$ is given by a closed subset of $Y\times Y$.
Consider the pullback square
\begin{center}
\begin{tikzcd}
	W & {\hom({\bbox}^{n+1}_\ud,X)} \\
	{\hom({\bbox}^{n+1}_\ud,X)} & {\hom({\bbcor}^{n+1}_\ud,X)}
	\arrow[from=1-1, to=1-2]
	\arrow[from=1-1, to=2-1]
	\arrow[from=1-2, to=2-2]
	\arrow[from=2-1, to=2-2]
\end{tikzcd}.
\end{center}
The condensed set $W$ is qcqs since all other vertices of the diagram are,
and we have that $Y\times_Z Y$ is the image of the map
\[
    W\to X(n+1)\times X(n+1)\to X(0)\times X(0) = Y\times Y
\]
by Proposition~\ref{prop:trunc_concfib}.
But images of qcqs objects in quasiseparated objects are qcqs again, so that $Y\times_Z Y$ is closed in $Y\times Y$.
\end{proof}

\begin{definition}
A condensed cubeset $X$ is \emph{$n$-ergodic} if $\tau_{n-1} X = \ast$.
\end{definition}

In particular, if $X$ is concrete and fibrant, then $X$ is $n$-ergodic if and only if $X(n) = X^{2^n}$.


\subsection{Structure Groups}

There also is the enrichment on structure groups.
For the following statement, we have to leave the setting of presentable categories and work externally in $\Cat$.

\begin{proposition}
Let $n\in\N$.
The functor $\pi_n\colon\cSet_\ast\to\cMon$ upgrades to a functor
\[
    \pi_n\colon\cCond_\ast\to\cMon(\CondSet) = \Cond(\cMon).
\]
We have that $\pi_n(X,x_0)(E) = \pi_n(X(E),x_0)$.
If $X$ is concrete fibrant, then $\pi_n(X,x_0)$ is a condensed abelian group.
If $X$ additionally is qcqs, so is $\pi_n(X,x_0)$.
\end{proposition}

\begin{proof}
The functor $\pi_n$ induces by pushforward a functor $\Fun(\extr_\kappa^\op,\cSet_\ast)\to\Fun(\extr_\kappa^\op,\cMon)$.
Since $\pi_n$ preserves finite products by Proposition~\ref{prop:structure-monoid-prod}
it restricts to a functor
\[
    \cSet_\ast\otimes\CondSet\simeq\Fun^\times(\extr_\kappa^\op,\cSet_\ast)\to\Fun^\times(\extr_\kappa^\op,\cMon)\simeq\cond(\cMon).
\]
If $X$ is qcqs, then $\tau_n X$ is as well, see Proposition~\ref{prop:cond-cube-trunc-qcqs},
and so the last claim follows from the fact that if $X$ is $n$-step, then $\pi_n(X,x_0)$ fits in the pullback square
\begin{center}
\begin{tikzcd}
	{\pi_n(X,x_0)} & \ast \\
	{\hom({\bbox}^n,X)} & {\hom({\bbcor}^n,X)}
	\arrow[from=1-1, to=1-2]
	\arrow[from=1-1, to=2-1]
	\arrow["{x_0}", from=1-2, to=2-2]
	\arrow[from=2-1, to=2-2]
\end{tikzcd}.
\end{center}
\end{proof}

\begin{remark}
The analogous statement holds for $\pi_0$:
since $\pi_0$ is left adjoint, tensoring with $\CondSet$ gives a functor $\pi_0\colon\cCond\to\CondSet$.
The functor is given by pointwise application of $\pi_0\colon\cSet\to\set$ since $\pi_0$ preserves finite products.
Moreover, it is computed by $\pi_0(X) = \tau_0 X(0)$.
\end{remark}

\begin{proposition}
Let $n\in\N$.
Then the functor $\cD_n\colon\Ab\to\cSet$ upgrades to a functor $\cD_n\colon\CondAb\to\cCond$
which factors through the subcategory of condensed abelian cubegroups.
We have that $\cD_n(A)(E) = \cD_n(A(E))$.
\end{proposition}

\begin{proof}
The pointwise pushforward $\cD_n\circ -\colon\Fun(\extr_\kappa^\op,\Ab)\to\Fun(\extr_\kappa^\op,\cSet)$
restricts to a functor $\cD_n\colon\CondAb\to\cCond$ since it preserves finite products by Remark~\ref{rem:EM_preserves_products}.
\end{proof}


\subsection{Weak Structure Theorem}

Now we are ready to state and prove a weak structure theorem for condensed cubesets.

\begin{theorem}\label{thm:cond-weak-structure}
Let $X$ be a condensed concrete ergodic fibrant cubeset.
Then there exists a tower of surjective fibrations
\begin{center}
\begin{tikzcd}
	X & \cdots & {\tau_2X} & {\tau_1X} & {\tau_0X}
	\arrow[from=1-1, to=1-2]
	\arrow[from=1-2, to=1-3]
	\arrow[from=1-3, to=1-4]
	\arrow[from=1-4, to=1-5]
\end{tikzcd}
\end{center}
such that $\tau_n X$ is a condensed concrete fibrant $n$-step cubeset
and $\tau_n X$ is a principal $\cD_n(\pi_n(X,x_0))$-bundle over $\tau_{n-1} X$.
This tower is natural in $X$.
\end{theorem}

Note that by projectivity of the point in $\cCond$, there is always a choice of some basepoint for nonempty $X$.

\begin{proof}
The theorem follows from the pointwise statement for the cubeset $X(E)$ from Theorem~\ref{thm:weak-structure}.
Note that the fact that the fibrations $\tau_n X\to\tau_{n-1} X$ are quotient maps
is automatic since in $\CondSet$, every epimorphism is regular.
Let us also remark that for $\tau_n X\to\tau_{n-1} X$ to be a principal $\cD_n(\pi_n(X,x_0))$-bundle,
it suffices to show that the shear map $\cD_n(\pi_n(X,x_0))\times\tau_n X\to\tau_n X\times_{\tau_{n-1}X}\tau_n X$ is an isomorphism,
since then the diagram
\begin{center}
\begin{tikzcd}
	{\cD_n(\pi_n(X,x_0))\times\tau_nX} & {\tau_n X} & {\tau_{n-1}X}
	\arrow[shift left, from=1-1, to=1-2]
	\arrow[shift right, from=1-1, to=1-2]
	\arrow[from=1-2, to=1-3]
\end{tikzcd}
\end{center}
is a coequalizer by surjectivity (= effective epimorphism) of $\tau_nX\to\tau_{n-1}X$.
\end{proof}

\begin{remark}
Recall that in this context, a principal $G$-bundle is a principal bundle in the wide sense,
i.e. without asserting local triviality of the action.
This is in line with \cite[Definition 2.1.6]{Candela2017a} and \cite[Theorem 5.4]{Gutman2020}
where no assumption on local triviality is made.
\end{remark}

By Proposition~\ref{prop:cond-cube-trunc-qcqs} the condensed weak structure theorem has the following corollary.

\begin{corollary}\label{cor:compact-weak-structure}
Let $X$ be a compact ergodic nilspace.
Then there is a tower of surjective fibrations
\begin{center}
\begin{tikzcd}
	X & \cdots & {\tau_2X} & {\tau_1X} & {\tau_0X}
	\arrow[from=1-1, to=1-2]
	\arrow[from=1-2, to=1-3]
	\arrow[from=1-3, to=1-4]
	\arrow[from=1-4, to=1-5]
\end{tikzcd}
\end{center}
such that $\tau_n X$ is a compact $n$-step nilspace
and the compact group $\pi_n(X,x_0)$ acts freely on $\tau_n X$
such that the orbits are the fibers of the quotient $\tau_n X\to\tau_{n-1} X$.
\end{corollary}

We have broken our convention and used classical terminology here,
the reason being an easier comparability to the classical theorem,
see \cite[Theorem 5.4]{Gutman2020} and \cite[Proposition 2.1.9]{Candela2017a}.
Note that a nilspace is an $n$-step concrete fibrant cubeset for some $n\in\N$,
see \cite[Definition 3.6]{Gutman2020} and \cite[Definition 1.2.1]{Candela2017}, \cite[Definition 1.0.2]{Candela2017a},
and as usual, $X$ is notationally identified with $X(0)$.
Moreover, compact is used to denote a compact second-countable Hausdorff space, see \cite[Definition 1.0.1]{Candela2017a}.

\begin{remark}
In general, there is no reason to believe that the principal bundle $\tau_n X\to\tau_{n-1}X$ is locally trivial,
i.e., the existence of a cover $U\to \tau_{n-1}X$ such that after base change to $U$,
the principal bundle is the trivial one $U\times G\to U$.
However, if $\pi_n(X,x_0)$ is a compact Lie group, then local triviality of the principal bundle is automatic
by \cite[Theorem 3.3]{Gleason1950}, see also \cite[Proposition 2.5.2]{Candela2017a}.
\end{remark}
